\section{Auto-correction : prédiction Dirichlet falsifiée}
\label{sec:auto_correction_dirichlet}
\label{sec:auto_correction}

La dualité K4 (\S\ref{sec:dualite_higgs_zeta}) suggère une extension
naturelle aux fonctions $L$ de Dirichlet primitives par factorisation
Haar $\bmod\,q$. Cette prédiction a été testée et falsifiée. La présente
section documente le test honnêtement et en tire la reformulation du
statut du résidu $1/8$.

\subsection{La prédiction P3}
\label{sec:P3}

Si K4 est structurellement extensible, la formule générique pour la
phase de Maslov de $L(s, \chi \bmod q)$ serait
\begin{equation}
\Delta_N(L \bmod q) \;=\; \sper^{2} \cdot \mathrm{Haar}(\Integers/q\Integers) \cdot \dfrac{1}{N_{\mathrm{branches}}} \;=\; \dfrac{1}{8q}.
\label{eq:P3}
\end{equation}
Prédictions~: $q = 3 : 1/24$, $q = 5 : 1/40$, $q = 7 : 1/56$,
$q = 11 : 1/88$.

\subsection{Test M1 (mpmath, $T \le 200$, 4 points par $q$)}
\label{sec:M1}

Premier test, inconclusif~: signal/bruit trop faible (4 points,
$\sigma_{\mathrm{SE}} \sim 0{,}2$, signal cible $\sim 0{,}04$).

\subsection{Test M1bis (LMFDB, $T \le 200$, ${\sim} 140$ zéros par $q$)}
\label{sec:M1bis}

Test approfondi sur les zéros pré-calculés de \texttt{LMFDB} pour les
caractères quadratiques primitifs $q \in \{3, 5, 7, 11\}$ ($a$ parité du
caractère)~:

\begin{table}[ht]
\centering
\small
\begin{tabular}{@{}rrrrrrrr@{}}
\toprule
$q$ & $a$ & $N$ zéros & $\langle \Delta_N \rangle$ & SE & $1/(8q)$ & $|t|_{\mathrm{PT}}$ & $|t|_{0}$\\
\midrule
3 & 1 & 114 & $-2{,}4 \cdot 10^{-4}$ & $0{,}019$ & $+0{,}042$ & \textbf{2{,}19} & $0{,}01$\\
5 & 0 & 129 & $-2{,}6 \cdot 10^{-3}$ & $0{,}021$ & $+0{,}025$ & $1{,}34$ & $0{,}13$\\
7 & 1 & 140 & $+1{,}8 \cdot 10^{-3}$ & $0{,}021$ & $+0{,}018$ & $0{,}77$ & $0{,}09$\\
11 & 1 & 155 & $-2{,}8 \cdot 10^{-4}$ & $0{,}021$ & $+0{,}011$ & $0{,}55$ & $0{,}01$\\
\bottomrule
\end{tabular}
\caption{Résultats du test M1bis. Le \emph{t-stat} compare la moyenne
observée à la prédiction PT $+1/(8q)$ (colonne $|t|_{\mathrm{PT}}$) et à
l'hypothèse nulle $\Delta_N = 0$ (colonne $|t|_{0}$).}
\label{tab:M1bis}
\end{table}

Test $\chi^{2}$ discriminatoire sur 4 d.d.l.~:

\begin{table}[ht]
\centering
\begin{tabular}{@{}lrr@{}}
\toprule
\textbf{Hypothèse} & $\chi^{2}$ & $p$-valeur\\
\midrule
$H_0$ : $\Delta_N = 0$ & $0{,}023$ & $\approx 1{,}00$\\
$H_{\mathrm{PT}}$ : $\Delta_N = +1/(8q)$ & $7{,}49$ & $0{,}11$\\
$H_{-\mathrm{PT}}$ : $\Delta_N = -1/(8q)$ & $7{,}05$ & $0{,}13$\\
\bottomrule
\end{tabular}
\caption{Test $\chi^{2}$ : $H_0$ est préféré par $\sim 300\times$ sur
$H_{\mathrm{PT}}$. La prédiction P3 est \textbf{effectivement falsifiée}
à $T \le 200$.}
\label{tab:chi2_dirichlet}
\end{table}

\paragraph{Discussion du signe.}
Sur les quatre valeurs de $q$ testées, PT prédit un décalage strictement
positif $+1/(8q)$ ; on observe trois décalages négatifs ($q = 3, 5, 11$)
et un positif ($q = 7$). L'accord en signe seul est donc $1/4$, bien
en-deçà du $1$ qu'imposerait la prédiction. La falsification n'est pas
seulement d'amplitude~: la structure de signe elle-même contredit la
prédiction $+1/(8q) > 0$.

\subsection{Calibration sur $\zeta$}
\label{sec:calibration_zeta}

Avec la même méthodologie sur les 50 premiers zéros de $\zeta$
(\cite{Odlyzko1992})~:
\begin{equation}
\langle \Delta_N \rangle_\zeta \;=\; +1{,}004 \pm 0{,}028.
\label{eq:calib_zeta}
\end{equation}
Comparaisons~:

\begin{table}[ht]
\centering
\begin{tabular}{@{}lrr@{}}
\toprule
\textbf{Cible} & \textbf{Valeur} & $t$\\
\midrule
Pôle de $\zeta$ seul ($+1$) & $1{,}000$ & $+0{,}15$ (compatible)\\
Pôle + Maslov ($1 + 1/8 = 1{,}125$) & $1{,}125$ & $-4{,}29$ (rejeté)\\
Maslov seul ($1/8$) & $0{,}125$ & $+31{,}29$ (rejeté)\\
\bottomrule
\end{tabular}
\caption{Calibration sur $\zeta$~: le $+1$ du pôle absorbe entièrement le
$1/8$ dans la statistique de comptage.}
\label{tab:calibration_zeta}
\end{table}

\subsection{Reformulation du statut du résidu $1/8$}
\label{sec:reformulation_residu}

\paragraph{Conséquence critique.}
La proposition~\ref{prop:densite_NPT} affirme que
$\NPT(\gamma) - N_{\mathrm{RvM}}(\gamma) = 1/8$, vérifié à $10^{-15}$ sur
quatre décades. Cette différence est entre \textbf{deux formes
asymptotiques lisses} ($\NPT$ et $N_{\mathrm{RvM}}$), pas un résidu
mesurable au comptage exact des zéros. Au comptage de zéros réels, le
$+1/8$ entre PT et RvM est absorbé par les autres termes de la formule
et n'apparaît pas comme un décalage moyen mesurable. Le $+1$ du pôle de
$\zeta$ absorbe entièrement le $1/8$ dans la statistique de comptage.

L'identité asymptotique~\eqref{eq:residu_un_huitieme} reste valide en
tant qu'identité \textbf{asymptotique entre formes}, mais la traduction
« $1/8$ visible dans le comptage des zéros » est erronée. Cela explique
a posteriori pourquoi la prédiction $1/(8q)$ a été falsifiée~: elle
attribuait à Dirichlet un résidu mesurable analogue à celui de $\zeta$,
alors que même pour $\zeta$ ce résidu n'est pas mesurable au comptage des
zéros à hauteur $T \sim 200$.

\subsection{Position dans le programme}
\label{sec:P3_position}

La falsification de P3 ne touche pas la structure HP-PT. Restent
intacts~:
\begin{itemize}[nosep,leftmargin=*]
\item le modèle spectral PT des zéros de $\zeta$ via $\Treg(s)$ :
$17/17$ zéros capturés à $|\Delta| < 1{,}5$ sur $t \in [10, 70]$,
$74/79$ sur $t \in [10, 200]$ avec la formule brute, $75/79$ à
$|\Delta| < 0{,}5$ après régularisation explicite $R_1$ avec
$\eps = 0{,}2$, et $267/269$ à $|\Delta| < 1{,}5$ étendu à
$t \in [10, 500]$ ;
\item l'identification cohomologique
$1/8 = \sper^{2}/2 = c_1/N_{\mathrm{corners}} = \sper/4$
pour $\zeta$ ;
\item la dualité Higgs $\leftrightarrow \zeta$ via le canal $p = 2$
(conjecture~\ref{conj:K4}) en tant que conjecture forte structurelle
pour $\zeta$.
\end{itemize}

Ce qui tombe~: l'extension naturelle à Dirichlet via la mesure de Haar
$\bmod\,q$. Le mécanisme PT décrit ici est \textbf{spécifique à $\zeta$},
non à la famille Dirichlet via cette factorisation. Le \emph{ledger} PT
marque P3 comme \textbf{[FALSIFIED, $T \le 200$]}.

C'est la discipline de falsifiabilité du programme~: une prédiction PT
testable a été produite par K4, testée, rejetée par les données, et le
programme se restructure proprement.

